\scanpage{134}
\mspage{83}

On sait que $q_{\dot T},q_T$ sont des iso, donc dans $\mathcal W$,
donc le ``lemme des cinq'' $\mathrm W(8)$ nous assure que
$\overline p\in\mathcal W\Longleftrightarrow q\unsure{,}\in\mathcal W$, qed.

\uline{Commentaire}. Cette proposition nous amène à
chercher des conditions sur un ouvert $U$ d'un
$X$, qui assurent que l'inclusion $i_0:U\to\overline U$
soit $\in\mathcal W$, et que \uline{cette propriété soit conservée
par tout cochangement de base} $U\unsure{-|\to}U'$. De
plus, nous recherchons des conditions qui soient
de nature locale sur $\overline U$, plus exactement sur
$\overline U\cap T=\dot U=\dot T$. Nous pouvons motiver ce
souci (a posteriori) en disant que nous
voudrions que la propriété $i_0\in\mathcal W$ soit conservée, non seulement par cochangement de base
\deleted{\unsure{que}}\inserted{quelconque} $U\to U'$, mais aussi, de plus, par
passage de $X$ à toute sous-catégorie pleine
$X_1$ contenant $U$, donnant donc naissance
à une situation
\[
 U\overset{\unsure{i_y|i_1}}{\subset} X_1,\qquad
 U\xrightarrow{(\unsure{d_y|i_y|i_1})_0}\overline U^{X_1}=\overline U\cap X_1
\]
\unsure{semblable|similaire} à la situation de départ. Plus
généralement encore, nous voudrions qu'elle
soit conservée en passant du \deleted{$(U,X)$}
triple $(X,U,T)$ à un triple ``induit'' $(X_1,U,T_1)$,
par \deleted{\unsure{la}} une \deleted{\unsure{fonction}}\inserted{flèche} quelconque $g:T_1\to T$ de
but $T$ dans $\mathrm{Cat}$, donnant naissance au
couple de flèches
\[
 U\xleftarrow{\mathrm{id}}U,\quad g:T_1\to T
\]
\unsure{cf pag|cf pp.} 60, 61). Le cas précédent est celui où

% archive sheet number: 134
